\begin{tikzpicture}[x=\spxunitstat,y=1.30cm]
\tikzset{
  bandTitle/.style={font=\fontsize{9.4}{10.3}\selectfont\sffamily\bfseries,
    text=VIUDark,anchor=west,inner sep=0pt},
  tag/.style={circle,draw=VIUDark,fill=VIUDark,text=white,
    minimum size=6.0mm,inner sep=0pt,
    font=\fontsize{8.1}{8.1}\selectfont\sffamily\bfseries},
  boxTitle/.style={font=\fontsize{7.9}{8.8}\selectfont\sffamily\bfseries,
    text=VIUDark,align=center},
  boxBody/.style={font=\fontsize{7.0}{7.8}\selectfont\sffamily,
    text=VIUGray,align=center},
  factorBody/.style={font=\fontsize{6.4}{7.1}\selectfont\sffamily,
    text=VIUGray,align=center},
  smallTitle/.style={font=\fontsize{7.6}{8.4}\selectfont\sffamily\bfseries,
    text=VIUDark,anchor=west,inner sep=0pt},
  smallBody/.style={font=\fontsize{7.0}{7.8}\selectfont\sffamily,
    text=VIUGray,anchor=west,inner sep=0pt},
  flow/.style={-{Latex[length=1.8mm,width=1.2mm]},draw=VIUOrange,
    line width=0.9pt},
  subflow/.style={-{Latex[length=1.5mm,width=1.0mm]},draw=VIUGray,
    line width=0.6pt},
  rule/.style={draw=SPGrid!88,line width=0.45pt}
}

\fill[white] (0.01,0.15) rectangle (13.49,14.10);
\draw[sparena] (0.01,0.15) rectangle (13.49,14.10);

% El título declara el alcance completo del protocolo.
\node[font=\fontsize{10.0}{10.9}\selectfont\sffamily\bfseries,
  text=VIUDark,anchor=west] at (0.40,13.72)
  {Protocolo Monte Carlo · campaña pareada};
\node[font=\fontsize{7.2}{8.0}\selectfont\sffamily,
  text=VIUOrange,anchor=east] at (13.08,13.72)
  {UNIDAD INDEPENDIENTE: MUNDO--SEMILLA};
\draw[draw=VIUDark,line width=0.85pt] (0.40,13.38) -- (13.08,13.38);

% A. Diseño factorial.
\node[tag] at (0.62,12.90) {A};
\node[bandTitle] at (1.00,12.90) {Diseño experimental y tamaño de la campaña};

\fill[SPBlue!3.0,rounded corners=2pt] (0.55,11.30) rectangle (3.28,12.48);
\draw[draw=SPGrid!90,rounded corners=2pt] (0.55,11.30) rectangle (3.28,12.48);
\node[boxTitle] at (1.915,12.22) {Escenarios $\mathcal C$};
\node[factorBody] at (1.915,11.68)
  {aleatorio · agrupado\\separado · anillo\\pasillo · fallo};

\fill[SPBlue!2.0,rounded corners=2pt] (3.48,11.30) rectangle (6.21,12.48);
\draw[draw=SPGrid!90,rounded corners=2pt] (3.48,11.30) rectangle (6.21,12.48);
\node[boxTitle] at (4.845,12.22) {Parámetros $\Theta$};
\node[factorBody] at (4.845,11.68)
  {$N,K$ · capacidad\\heterogeneidad · red\\perturbación · timeout};

\fill[SPGreen!2.5,rounded corners=2pt] (6.41,11.30) rectangle (9.14,12.48);
\draw[draw=SPGreen!55,rounded corners=2pt] (6.41,11.30) rectangle (9.14,12.48);
\node[boxTitle] at (7.775,12.22) {Semillas $\mathcal S$};
\node[font=\fontsize{9.5}{10.3}\selectfont\sffamily,text=SPGreen!55!black]
  at (7.775,11.72) {$\{s_1,\ldots,s_S\}$};

\fill[VIUOrange!2.5,rounded corners=2pt] (9.34,11.30) rectangle (12.95,12.48);
\draw[draw=VIUOrange!55,rounded corners=2pt] (9.34,11.30) rectangle (12.95,12.48);
\node[boxTitle] at (11.145,12.22) {Métodos $\mathcal M$};
\node[factorBody] at (11.145,11.66)
  {propuesta · baselines\\oráculo cuando sea certificable};

% Producto correcto: primero mundos; después métodos sobre el mismo mundo.
\draw[flow] (1.915,11.24) -- (1.915,11.08);
\draw[flow] (4.845,11.24) -- (4.845,11.08);
\draw[flow] (7.775,11.24) -- (7.775,11.08);
\fill[SPGreen!2.2,rounded corners=2pt] (0.78,10.35) rectangle (8.78,10.98);
\draw[draw=SPGreen!55,rounded corners=2pt] (0.78,10.35) rectangle (8.78,10.98);
\node[font=\fontsize{8.6}{9.4}\selectfont\sffamily\bfseries,
  text=VIUDark] at (4.78,10.665)
  {mundos $\mathcal W=\mathcal C\times\Theta\times\mathcal S$
   \quad $|\mathcal W|=|\mathcal C|\,|\Theta|\,S$};
\draw[flow] (8.84,10.665) -- (9.22,10.665);
\draw[flow] (11.145,11.24) -- (11.145,11.08);
\fill[VIUOrange!2.2,rounded corners=2pt] (9.28,10.35) rectangle (12.95,10.98);
\draw[draw=VIUOrange!55,rounded corners=2pt] (9.28,10.35) rectangle (12.95,10.98);
\node[font=\fontsize{8.1}{8.9}\selectfont\sffamily\bfseries,
  text=VIUDark] at (11.115,10.665)
  {ejecuciones $|\mathcal R|=|\mathcal W|\,|\mathcal M|$};
\draw[rule] (0.40,10.22) -- (13.08,10.22);

% B. Bucle de ejecución Monte Carlo como diagrama de control.
\node[tag] at (0.62,9.91) {B};
\node[bandTitle] at (1.00,9.91) {Bucle Monte Carlo pareado y registro RAW};

% Una iteración procesa un mundo; el bus pareado ejecuta todos los métodos.
\fill[SPBlue!1.3,rounded corners=2pt] (0.55,7.55) rectangle (12.95,9.48);
\draw[draw=SPGrid!82,rounded corners=2pt,line width=0.5pt]
  (0.55,7.55) rectangle (12.95,9.48);

\fill[white,rounded corners=2pt] (0.72,8.16) rectangle (2.38,9.25);
\draw[draw=SPBlue!48,rounded corners=2pt] (0.72,8.16) rectangle (2.38,9.25);
\node[font=\fontsize{6.8}{7.5}\selectfont\sffamily\bfseries,text=VIUDark]
  at (1.55,9.02) {1 · CONTADOR};
\node[font=\fontsize{6.5}{7.2}\selectfont\sffamily\bfseries,text=VIUDark]
  at (1.55,8.64) {$r=1,\ldots,|\mathcal W|$};
\node[font=\fontsize{6.4}{7.1}\selectfont\sffamily,text=VIUGray]
  at (1.55,8.35) {$(c_r,\theta_r,s_r)$};

\draw[flow] (2.42,8.70) -- (2.66,8.70);
\fill[SPGreen!3.0,rounded corners=2pt] (2.72,8.16) rectangle (4.55,9.25);
\draw[draw=SPGreen!58,rounded corners=2pt] (2.72,8.16) rectangle (4.55,9.25);
\node[font=\fontsize{6.8}{7.5}\selectfont\sffamily\bfseries,text=VIUDark]
  at (3.635,9.02) {2 · GENERADOR};
\node[font=\fontsize{7.4}{8.2}\selectfont\sffamily\bfseries,text=VIUDark]
  at (3.635,8.64) {$w_r=G(\xi_r)$};
\node[font=\fontsize{6.1}{6.8}\selectfont\sffamily,text=VIUGray]
  at (3.635,8.35) {$\xi_r=(c_r,\theta_r,s_r)$};

\draw[flow] (4.59,8.70) -- (4.84,8.70);
\fill[white,rounded corners=2pt] (4.90,8.02) rectangle (7.43,9.38);
\draw[draw=VIUOrange!58,rounded corners=2pt] (4.90,8.02) rectangle (7.43,9.38);
\node[font=\fontsize{7.0}{7.7}\selectfont\sffamily\bfseries,text=VIUDark]
  at (6.165,9.15) {3 · MÉTODOS};
% Fan-out y fan-in: la geometría hace visible el mismo mundo compartido.
\draw[draw=VIUOrange,line width=0.75pt] (5.13,8.23) -- (5.13,8.75);
\draw[draw=VIUOrange,line width=0.75pt] (5.13,8.49) -- (5.46,8.49);
\foreach \yy/\lab in {8.75/$m_1(w_r)$,8.49/$\vdots$,8.23/$m_{|\mathcal M|}(w_r)$}{
  \draw[draw=VIUOrange,line width=0.65pt] (5.13,\yy) -- (5.46,\yy);
  \fill[VIUOrange!3,rounded corners=1pt] (5.46,\yy-0.105) rectangle (6.78,\yy+0.105);
  \draw[draw=VIUOrange!48,rounded corners=1pt] (5.46,\yy-0.105) rectangle (6.78,\yy+0.105);
  \node[font=\fontsize{6.4}{7.0}\selectfont\sffamily,text=VIUDark] at (6.12,\yy) {\lab};
  \draw[draw=VIUOrange,line width=0.65pt] (6.78,\yy) -- (7.10,\yy);
}
\draw[draw=VIUOrange,line width=0.75pt] (7.10,8.23) -- (7.10,8.75);
\draw[flow] (7.47,8.70) -- (7.67,8.70);
\fill[SPFloor!65,rounded corners=2pt] (7.73,8.16) rectangle (9.53,9.25);
\draw[draw=SPGrid!96,rounded corners=2pt] (7.73,8.16) rectangle (9.53,9.25);
\node[font=\fontsize{7.0}{7.7}\selectfont\sffamily\bfseries,text=VIUDark]
  at (8.63,9.02) {4 · REGISTRO};
\node[font=\fontsize{7.8}{8.6}\selectfont\sffamily\bfseries,text=VIUDark]
  at (8.63,8.64) {guardar RAW};
\node[font=\fontsize{6.4}{7.1}\selectfont\sffamily,text=VIUGray]
  at (8.63,8.35) {$m$ · estado · $Y_{w_r}^m$};

\draw[flow] (9.59,8.70) -- (9.83,8.70);
\node[diamond,draw=VIUDark!72,fill=white,aspect=1.75,
  minimum width=14.8mm,minimum height=8.8mm,inner sep=0.8pt,
  font=\fontsize{6.6}{7.3}\selectfont\sffamily\bfseries,
  text=VIUDark,align=center] (loopgate) at (10.52,8.70)
  {$r<|\mathcal W|$?};

\draw[flow] (loopgate.east) -- (11.49,8.70);
\node[font=\fontsize{6.2}{6.9}\selectfont\sffamily\bfseries,
  text=VIUOrange,fill=SPBlue!1.3,inner sep=1pt] at (11.32,8.94) {NO};
\fill[SPGreen!3,rounded corners=6pt] (11.55,8.27) rectangle (12.90,9.13);
\draw[draw=SPGreen!60,rounded corners=6pt] (11.55,8.27) rectangle (12.90,9.13);
\node[font=\fontsize{7.0}{7.7}\selectfont\sffamily\bfseries,
  text=SPGreen!45!black] at (12.225,8.86) {FIN};
\node[font=\fontsize{6.2}{6.9}\selectfont\sffamily\bfseries,
  text=SPGreen!45!black] at (12.225,8.54) {RAW cerrado};

% Realimentación del controlador para la siguiente unidad mundo--semilla.
\draw[draw=VIUOrange,line width=0.85pt,-{Latex[length=1.7mm]}]
  (loopgate.south) -- (10.52,7.72) -- (1.55,7.72) -- (1.55,8.08);
\node[font=\fontsize{6.6}{7.3}\selectfont\sffamily\bfseries,
  text=VIUOrange,fill=SPBlue!1.3,inner sep=1.4pt] at (6.10,7.72)
  {SÍ · $r\leftarrow r+1$};

% Esquema mínimo que hace auditable cada fila sin convertirlo en tabla.
\fill[SPFloor!58,rounded corners=2pt] (0.72,6.72) rectangle (12.82,7.36);
\draw[draw=SPGrid!92,rounded corners=2pt] (0.72,6.72) rectangle (12.82,7.36);
\node[font=\fontsize{6.8}{7.5}\selectfont\sffamily\bfseries,
  text=VIUDark,anchor=west] at (0.95,7.04) {RAW INMUTABLE};
\draw[draw=SPGrid!88,line width=0.45pt] (3.05,6.82) -- (3.05,7.26);
\node[font=\fontsize{6.2}{6.9}\selectfont\sffamily,text=VIUGray,
  anchor=west] at (3.25,7.04)
  {ID · $c$ · $\theta$ · $s$ · $m$ · versión · estado · $Y$ · CPU · certificado};
\draw[draw=SPGrid!88,line width=0.45pt] (10.42,6.82) -- (10.42,7.26);
\node[font=\fontsize{6.2}{6.9}\selectfont\sffamily\bfseries,
  text=VIUOrange,anchor=east] at (12.60,7.04)
  {FALLO/TIMEOUT: visibles};
\draw[rule] (0.40,6.50) -- (13.08,6.50);

% C. Procesamiento reproducible.
\node[tag] at (0.62,6.11) {C};
\node[bandTitle] at (1.00,6.11) {Procesamiento reproducible: RAW $\rightarrow$ evidencia por mundo};

\fill[SPFloor!62,rounded corners=2pt] (0.55,4.58) rectangle (2.55,5.72);
\draw[draw=SPGrid!92,rounded corners=2pt] (0.55,4.58) rectangle (2.55,5.72);
\node[boxTitle] at (1.55,5.45) {1 · RAW};
\node[font=\fontsize{6.6}{7.4}\selectfont\sffamily,text=VIUGray,align=center]
  at (1.55,4.95) {fila por $m(w)$\\sin descartes};

\draw[flow] (2.62,5.15) -- (2.87,5.15);
\fill[SPBlue!2.5,rounded corners=2pt] (2.94,4.58) rectangle (6.12,5.72);
\draw[draw=SPBlue!45,rounded corners=2pt] (2.94,4.58) rectangle (6.12,5.72);
\node[boxTitle] at (4.53,5.45) {2 · MÉTRICAS};
\node[font=\fontsize{7.7}{8.5}\selectfont\sffamily\bfseries,text=VIUDark]
  at (4.53,4.95) {$Y_w^m=(F,J,t,\mathrm{CPU},\ldots)$};

\draw[flow] (6.19,5.15) -- (6.48,5.15);
\fill[VIUOrange!2.2,rounded corners=2pt] (6.55,4.58) rectangle (9.54,5.72);
\draw[draw=VIUOrange!55,rounded corners=2pt] (6.55,4.58) rectangle (9.54,5.72);
\node[boxTitle] at (8.045,5.45) {3 · PAREAMIENTO};
\node[font=\fontsize{9.0}{9.8}\selectfont\sffamily\bfseries,text=VIUOrange]
  at (8.045,4.96) {$\Delta_w^{A,B}=Y_w^A-Y_w^B$};

\draw[flow] (9.61,5.15) -- (9.86,5.15);
\fill[SPGreen!2.2,rounded corners=2pt] (9.93,4.58) rectangle (12.95,5.72);
\draw[draw=SPGreen!55,rounded corners=2pt] (9.93,4.58) rectangle (12.95,5.72);
\node[boxTitle] at (11.44,5.45) {4 · TABLA ANALÍTICA};
\node[boxBody] at (11.44,4.95) {fila por mundo $w$\\respuesta · estado · $\Delta$};
\node[font=\fontsize{7.0}{7.8}\selectfont\sffamily\bfseries,
  text=SPGreen!55!black,align=center] at (6.75,4.27)
  {bootstrap y contrastes remuestrean/comparan mundos; robots y cargas permanecen anidados};
\draw[rule] (0.40,3.96) -- (13.08,3.96);

% D. Inferencia y decisión sobre hipótesis.
\node[tag] at (0.62,3.58) {D};
\node[bandTitle] at (1.00,3.58) {Inferencia y conclusión sobre la hipótesis preespecificada};
\draw[rule] (5.05,1.45) -- (5.05,3.30);
\draw[rule] (8.05,1.45) -- (8.05,3.30);

\node[smallTitle] at (0.62,3.14) {CONTRASTE SEGÚN RESPUESTA};
\node[smallBody] at (0.62,2.75) {binaria $\rightarrow$ McNemar exacto};
\node[smallBody] at (0.62,2.41) {continua $\rightarrow$ Wilcoxon · $r_{rb}$};
\node[smallBody] at (0.62,2.07) {$\geq3$ métodos $\rightarrow$ Friedman · Kendall $W$};
\node[smallBody] at (0.62,1.73) {censura $\rightarrow$ supervivencia / $t$ truncado};

\node[smallTitle] at (5.35,3.14) {ESTIMACIÓN};
\node[font=\fontsize{9.7}{10.6}\selectfont\sffamily\bfseries,
  text=VIUDark] at (6.55,2.63)
  {$\widehat\Delta,\ \mathrm{IC}_{95\%},\ p_{\mathrm{Holm}}$};
\node[smallBody] at (5.35,2.17) {efecto · tasa de fallo};
\node[smallBody] at (5.35,1.83) {gap: oráculo certificado};

\node[smallTitle] at (8.34,3.14) {DECISIÓN $H_0/H_1$};
\node[font=\fontsize{6.7}{7.4}\selectfont\sffamily,
  text=VIUGray,anchor=west] at (8.34,2.82)
  {PRE: fijar $H_0$, $H_1$, $\alpha$ y umbral práctico};
\fill[SPFloor!70,rounded corners=2pt] (8.34,2.28) rectangle (12.90,2.64);
\node[font=\fontsize{6.8}{7.5}\selectfont\sffamily\bfseries,
  text=VIUDark] at (10.62,2.46)
  {$p_{Holm}\leq\alpha$ y ¿efecto/IC cumplen el umbral?};
\draw[subflow] (10.62,2.26) -- (9.49,2.08);
\draw[subflow] (10.62,2.26) -- (11.77,2.08);
\fill[SPGreen!5,rounded corners=2pt] (8.34,1.46) rectangle (10.58,2.02);
\draw[draw=SPGreen!60,rounded corners=2pt] (8.34,1.46) rectangle (10.58,2.02);
\node[font=\fontsize{6.7}{7.4}\selectfont\sffamily\bfseries,
  text=SPGreen!45!black,align=center] at (9.46,1.74)
  {SÍ: rechazar $H_0$\\$H_1$ sustentada};
\fill[VIUOrange!4,rounded corners=2pt] (10.70,1.46) rectangle (12.90,2.02);
\draw[draw=VIUOrange!60,rounded corners=2pt] (10.70,1.46) rectangle (12.90,2.02);
\node[font=\fontsize{6.6}{7.3}\selectfont\sffamily\bfseries,
  text=VIUOrange!82!black,align=center] at (11.80,1.74)
  {NO: no rechazar $H_0$\\$H_1$ no sustentada};

% Cierre: la decisión incluye evidencia y limitación, no solo p-valor.
\fill[VIUDark,rounded corners=2pt] (0.40,0.48) rectangle (13.08,1.25);
\node[font=\fontsize{8.3}{9.1}\selectfont\sffamily\bfseries,
  text=white,align=center] at (6.74,0.87)
  {CONCLUSIÓN: efecto + incertidumbre + decisión $H_0/H_1$ + limitaciones};
\end{tikzpicture}
